\chapter{Estrogen Receptor and Progesterone Receptor Tests}
\section*{What Is This Test?} These pathology tests detect estrogen and progesterone receptors in breast-cancer tissue and help predict response to hormone therapy.
\section*{Alternative Names} ER testing; PR testing; hormone-receptor testing; breast biomarker testing.
\section*{Principle} Immunohistochemistry stains tumor cells and reports the proportion and intensity of receptor expression using laboratory criteria.
\section*{Why This Test Is Ordered} It is performed on breast-cancer tissue to classify disease and guide systemic treatment decisions.
\section*{Primary vs Secondary Test} Receptor testing is a secondary pathology test after tissue diagnosis; it does not screen healthy people for breast cancer.
\section*{How This Test Is Performed} A pathologist tests tissue from core biopsy or surgery, with controls and quality review.
\section*{What Preparation Is Needed} No patient preparation beyond the biopsy or surgery is required.
\section*{Understanding Results} Positive, negative, or borderline results are interpreted with tumor type, grade, HER2, and clinical context. Discordant or weak results may need review or repeat testing.
\section*{Follow-up Tests} HER2 testing, genomic assays, staging studies, genetic counseling, and oncology consultation may follow.
\section*{References} \begin{enumerate}\item MedlinePlus. Estrogen Receptor, Progesterone Receptor Tests.\item American Society of Clinical Oncology and College of American Pathologists. Breast biomarker guideline.\end{enumerate}
\clearpage\section*{Understanding Results and Follow-up}
The laboratory report should identify the specimen, method, units, reference interval or decision limit, and whether the result is positive, negative, abnormal, or inconclusive. Reference intervals vary with age, sex, pregnancy, specimen type, assay, and laboratory; a result outside a range is not automatically a diagnosis, and a result within a range does not always exclude disease. Discuss unexpected results with the ordering clinician, who can decide whether repeat or confirmatory testing, treatment, or referral is appropriate. Do not change medicines or delay urgent care based on an isolated result.
\section*{Regulatory Status and Authorization Date}This chapter describes a test category, not a specific commercial product. FDA, CDSCO, EU IVDR, and MHRA approval or registration dates therefore cannot be assigned without the assay manufacturer, product name, instrument, and intended use. For laboratory-developed tests, an individual agency approval date may not exist. See the Preface for the interpretation of regulatory status.
\clearpage
